===== BEGIN FILE: main.tex =====
\documentclass[
 reprint,
superscriptaddress,









 amsmath,amssymb,
 aps,
prl,





]{revtex4-2}

\usepackage{graphicx}
\usepackage{dcolumn}
\usepackage{bm}











\usepackage{comment}
\usepackage[utf8]{inputenc}
\usepackage{soul}
\usepackage{xcolor}

\usepackage{booktabs}

\newcommand{\SY}[1]{\textcolor{teal}{#1}}
\newcommand{\SA}[1]{\textcolor{blue}{#1}}
\newcommand{\TBC}[1]{\textcolor{red}{#1}}


\begin{document}

\preprint{APS/123-QED}

\title{Extended Haloscope Search and Candidate Validation near 1.036\,GHz}





\author{Saebyeok Ahn}
\affiliation{Dark Matter Axion Group, Institute for Basic Science, Daejeon 34126, Republic of Korea}
\affiliation{Center for Axion and Precision Physics Research, Institute for Basic Science (IBS), Daejeon 34051, Republic of Korea}




\author{Boris I. Ivanov}
\affiliation{Dark Matter Axion Group, Institute for Basic Science, Daejeon 34126, Republic of Korea}
\affiliation{Center for Axion and Precision Physics Research, Institute for Basic Science (IBS), Daejeon 34051, Republic of Korea}

\author{Ohjoon Kwon}
\affiliation{Dark Matter Axion Group, Institute for Basic Science, Daejeon 34126, Republic of Korea}
\affiliation{Center for Axion and Precision Physics Research, Institute for Basic Science (IBS), Daejeon 34051, Republic of Korea}

\author{HeeSu Byun}
\affiliation{Center for Axion and Precision Physics Research, Institute for Basic Science (IBS), Daejeon 34051, Republic of Korea}

\author{Arjan F. van Loo}
\affiliation{RIKEN Center for Quantum Computing (RQC), Wako, Saitama 351-0198, Japan}
\affiliation{Department of Applied Physics, Graduate School of Engineering, The University of Tokyo, Bunkyo-ku, Tokyo 113-8656, Japan}

\author{SeongTae Park}
\affiliation{Center for Axion and Precision Physics Research, Institute for Basic Science (IBS), Daejeon 34051, Republic of Korea}

\author{JinMyeong Kim}
\affiliation{Department of Physics, Korea Advanced Institute of Science and Technology (KAIST), Daejeon 34141, Republic of Korea}
\affiliation{Center for Axion and Precision Physics Research, Institute for Basic Science (IBS), Daejeon 34051, Republic of Korea}

\author{Junu Jeong}
\thanks{Present address: Oskar Klein Centre, Department of Physics, Stockholm University, AlbaNova, SE-10691 Stockholm, Sweden}
\affiliation{Center for Axion and Precision Physics Research, Institute for Basic Science (IBS), Daejeon 34051, Republic of Korea}

\author{Soohyung Lee}
\thanks{Present address: Center for Accelerator Research, Korea University, Sejong 30019, Republic of Korea}
\affiliation{Center for Axion and Precision Physics Research, Institute for Basic Science (IBS), Daejeon 34051, Republic of Korea}

\author{Jinsu Kim}
\affiliation{Dark Matter Axion Group, Institute for Basic Science, Daejeon 34126, Republic of Korea}
\affiliation{Center for Axion and Precision Physics Research, Institute for Basic Science (IBS), Daejeon 34051, Republic of Korea}

\author{Çağlar Kutlu}
\thanks{Present address: Zurich Instruments, Technoparkstrasse 1, 8005 Zürich, Switzerland.}
\affiliation{Center for Axion and Precision Physics Research, Institute for Basic Science (IBS), Daejeon 34051, Republic of Korea}

\author{Andrew K. Yi}
\thanks{Present address: SLAC National Accelerator Laboratory, 2575 Sand Hill Road, Menlo Park, California 94025, USA}
\affiliation{Department of Physics, Korea Advanced Institute of Science and Technology (KAIST), Daejeon 34141, Republic of Korea}
\affiliation{Center for Axion and Precision Physics Research, Institute for Basic Science (IBS), Daejeon 34051, Republic of Korea}

\author{Yasunobu Nakamura}
\affiliation{RIKEN Center for Quantum Computing (RQC), Wako, Saitama 351-0198, Japan}
\affiliation{Department of Applied Physics, Graduate School of Engineering, The University of Tokyo, Bunkyo-ku, Tokyo 113-8656, Japan}

\author{Seonjeong Oh}
\affiliation{Dark Matter Axion Group, Institute for Basic Science, Daejeon 34126, Republic of Korea}
\affiliation{Center for Axion and Precision Physics Research, Institute for Basic Science (IBS), Daejeon 34051, Republic of Korea}

\author{Danho Ahn}
\thanks{Present address: INFN-Sezione di Padova, Via Marzolo 8, 35131 Padova, Italy}
\affiliation{Center for Axion and Precision Physics Research, Institute for Basic Science (IBS), Daejeon 34051, Republic of Korea}

\author{SungJae Bae}
\thanks{Present address: RIKEN Center for Quantum Computing (RQC), Wako, Saitama 351–0198, Japan}
\affiliation{Department of Physics, Korea Advanced Institute of Science and Technology (KAIST), Daejeon 34141, Republic of Korea}
\affiliation{Center for Axion and Precision Physics Research, Institute for Basic Science (IBS), Daejeon 34051, Republic of Korea}

\author{Hyoungsoon Choi}
\affiliation{Department of Physics, Korea Advanced Institute of Science and Technology (KAIST), Daejeon 34141, Republic of Korea}

\author{Jihoon Choi}
\thanks{Present address: Korea Astronomy and Space Science Institute, Daejeon 34055, Republic of Korea.}
\affiliation{Center for Axion and Precision Physics Research, Institute for Basic Science (IBS), Daejeon 34051, Republic of Korea}

\author{Yonuk Chong}
\affiliation{SKKU Advancd Institute of Nano Technology (SAINT) and Department of Nano Engineering
Sung Kyun Kwan University (SKKU), Suwon 16419, Republic of Korea}

\author{Woohyun Chung}
\affiliation{Center for Axion and Precision Physics Research, Institute for Basic Science (IBS), Daejeon 34051, Republic of Korea}

\author{Violeta Gkika}
\affiliation{Center for Axion and Precision Physics Research, Institute for Basic Science (IBS), Daejeon 34051, Republic of Korea}

\author{Jihn E. Kim}
\affiliation{Department of Physics, Seoul National University, 1 Gwanak-Ro, Seoul 08826, Republic of Korea}


\author{Younggeun Kim}
\thanks{Present address: Johannes Gutenberg-Universität Mainz, 55122 Mainz, Germany; GSI Helmholtzzentrum für Schwerionenforschung GmbH, 64291 Darmstadt, Germany}
\affiliation{Center for Axion and Precision Physics Research, Institute for Basic Science (IBS), Daejeon 34051, Republic of Korea}

\author{Byeong Rok Ko}
\thanks{Present address: Department of Accelerator Science, Korea University Sejong Campus, 2511 Sejong-ro, Sejong, 30019, Republic of Korea}
\affiliation{Center for Axion and Precision Physics Research, Institute for Basic Science (IBS), Daejeon 34051, Republic of Korea}

\author{Lino Miceli}
\affiliation{Center for Axion and Precision Physics Research, Institute for Basic Science (IBS), Daejeon 34051, Republic of Korea}

\author{Doyu Lee}
\thanks{Present address: Samsung Electronics, Gyeonggi-do 16677, Republic of Korea.}
\affiliation{Center for Axion and Precision Physics Research, Institute for Basic Science (IBS), Daejeon 34051, Republic of Korea}

\author{Jiwon Lee}
\affiliation{Dark Matter Axion Group, Institute for Basic Science, Daejeon 34126, Republic of Korea}
\affiliation{Department of Physics, Korea Advanced Institute of Science and Technology (KAIST), Daejeon 34141, Republic of Korea}
\affiliation{Center for Axion and Precision Physics Research, Institute for Basic Science (IBS), Daejeon 34051, Republic of Korea}

\author{Ki Woong Lee}
\affiliation{Center for Axion and Precision Physics Research, Institute for Basic Science (IBS), Daejeon 34051, Republic of Korea}

\author{MyeongJae Lee}
\thanks{Present address: Department of Physics, Sungkyunkwan University, Suwon 16419, Republic of Korea.}
\affiliation{Center for Axion and Precision Physics Research, Institute for Basic Science (IBS), Daejeon 34051, Republic of Korea}

\author{Andrei Matlashov}
\thanks{Deceased.}
\affiliation{Center for Axion and Precision Physics Research, Institute for Basic Science (IBS), Daejeon 34051, Republic of Korea}

\author{Pallavi Parashar}
\affiliation{Dark Matter Axion Group, Institute for Basic Science, Daejeon 34126, Republic of Korea}
\affiliation{Department of Physics, Korea Advanced Institute of Science and Technology (KAIST), Daejeon 34141, Republic of Korea}
\affiliation{Center for Axion and Precision Physics Research, Institute for Basic Science (IBS), Daejeon 34051, Republic of Korea}

\author{Taehyeon Seong}
\affiliation{Dark Matter Axion Group, Institute for Basic Science, Daejeon 34126, Republic of Korea}
\affiliation{Center for Axion and Precision Physics Research, Institute for Basic Science (IBS), Daejeon 34051, Republic of Korea}

\author{Yun Chang Shin}
\affiliation{Center for Axion and Precision Physics Research, Institute for Basic Science (IBS), Daejeon 34051, Republic of Korea}


\author{Sergey V. Uchaikin}
\affiliation{Dark Matter Axion Group, Institute for Basic Science, Daejeon 34126, Republic of Korea}
\affiliation{Center for Axion and Precision Physics Research, Institute for Basic Science (IBS), Daejeon 34051, Republic of Korea}



\author{Yannis K. Semertzidis}
\affiliation{Center for Axion and Precision Physics Research, Institute for Basic Science (IBS), Daejeon 34051, Republic of Korea}
\affiliation{Department of Physics, Korea Advanced Institute of Science and Technology (KAIST), Daejeon 34141, Republic of Korea}

\author{SungWoo Youn}
\thanks{corresponding author}
\email{swyoun@ibs.re.kr}
\affiliation{Dark Matter Axion Group, Institute for Basic Science, Daejeon 34126, Republic of Korea}
\affiliation{Center for Axion and Precision Physics Research, Institute for Basic Science (IBS), Daejeon 34051, Republic of Korea}


\date{\today}

\begin{abstract}
We report a follow-up axion haloscope search near 1.036 GHz that completes and extends our previous work [Phys. Rev. X \textbf{14}, 031023 (2024)], in which a portion of the HEMT-based data could not be analyzed due to unrecorded experimental information.
While recovering this dataset, we identified an excess near 1.036 GHz that satisfied our candidate-selection criteria, motivating dedicated validation studies, including independent cross-checks and re-examination with the original apparatus.
The excess did not persist under these investigations and was not confirmed as an axion dark-matter signal.
We subsequently extended the search over a 20-MHz band surrounding the candidate using a quantum-noise-limited amplifier, achieving sensitivity close to the Dine–Fischler–Srednicki–Zhitnitsky benchmark.
In the absence of a confirmed signal, we set improved 90\% confidence-level upper limits on the axion–photon coupling over the frequency range 1.026–1.045 GHz.
This work highlights the importance of robust candidate-validation strategies as haloscope searches approach discovery-level sensitivity.
\end{abstract}



\maketitle


The axion, originally proposed as a solution to the strong charge–parity problem in quantum chromodynamics~\cite{PRL37_8_1976,PRD14_3432_1978,PRD18_2199_1978,PR108_120_1957,PRD15_9_1977,NPA341_269_1980,PRL38_1440_1977,PRL40_223_1978,PRL40_279_1978}, is a well-motivated candidate for cold dark matter~\cite{AA594_A13_2016}.
Resonant microwave haloscopes search for axion dark matter through axion-to-photon conversion in a strong magnetic field, with the conversion power enhanced by a high-$Q$ cavity.
The expected signal power is given by~\cite{CAPP_scan_rate}
\begin{equation}
P_{\rm sig} = g_{a\gamma\gamma}^2 \, \frac{\rho_a}{m_a} \, B_0^2 V C Q \frac{\beta}{\beta+1},
\end{equation}
where $g_{a\gamma\gamma}$ denotes the axion–photon coupling, $\rho_a$ the local axion dark-matter density, $m_a$ the axion mass, $B_0$ the applied magnetic field strength, $V$ the cavity volume, $C$ the mode-dependent form factor, $Q$ the cavity quality factor, and $\beta$ the antenna coupling.
Benchmark axion models such as Kim–Shifman–Vainshtein–Zakharov~(KSVZ)~\cite{PRL43_103_1979,NPB166_493_1980} and Dine–Fischler–Srednicki–Zhitnitsky~(DFSZ)~\cite{YF31_497_1980,PLB104_199_1981} define the natural sensitivity goals for haloscope experiments.
As advances in cryogenic amplification bring receiver noise toward the quantum limit, haloscope experiments are entering a regime of discovery-level sensitivity.
In this regime, rigorous candidate identification and validation become essential to distinguish genuine axion signals from instrumental artifacts and environmental radio-frequency interference.

In a recent search reported in Ref.~\cite{CAPP-MAX}, we set stringent limits on the axion–photon coupling over a broad frequency range.
Most of the range was scanned using Josephson parametric amplifiers (JPAs)~\cite{CAPP_jpa_frontier,CAPP_jpa_char}, while certain intervals outside the JPA bandwidth were covered with high-electron-mobility transistor (HEMT) amplification.
A subset of data in the 1.033--1.037\,GHz range could not be analyzed due to missing antenna-coupling measurements and was therefore excluded from the original report.

In this Letter, we complete the missing frequency band to provide continuous exclusion coverage and establish a reference point for subsequent higher-sensitivity searches.
The restored data revealed a feature near 1.036\,GHz that met our candidate-selection criteria, prompting a comprehensive program of validation studies.
We report the outcome of these investigations: independent cross-checks with a separate haloscope system and re-examination using the original apparatus.
We subsequently extended the search using a JPA receiver and report updated 90\% confidence-level limits on the axion–photon coupling over the frequency range 1.026--1.045 GHz.

The original scan~\cite{CAPP-MAX} was performed with a haloscope, featuring a 12-T superconducting solenoid, a 37-L cylindrical copper cavity and near-quantum-limited JPAs.
The system was maintained below 40\,mK by a dilution refrigerator, substantially suppressing thermal noise from both the cavity and the receiver chain.
The experiment operated in the TM$_{010}$ cavity mode, which maximizes axion–photon conversion power, with frequency tuning achieved via a piezoelectric-actuated copper rod inside the cavity.
The cavity signal was extracted via a strongly coupled coaxial antenna with $\beta\sim1.5$.
The receiver chain comprised JPAs operated with gains of $\sim$20\,dB with a bandwidth of $\sim$100\,kHz, followed by cryogenic HEMT amplifiers and subsequent room-temperature amplification.
The cavity resonant frequency was set 100-kHz away from the JPA resonance in order to operate the JPA in the phase-insensitive mode.
The amplified signals were down-converted to an intermediate frequency of 10.7\,MHz using an image-rejection mixer and digitized with a resolution bandwidth of 10\,Hz.
A dedicated aerial antenna system was employed to monitor environmental radio-frequency interference.
A schematic of the experimental setup is shown in Fig.~\ref{fig:schematics}.

\begin{figure}[ht]
    \centering
    \includegraphics[width=0.9\linewidth]{figure/schematics.pdf}
    \caption{Schematic of the axion haloscope setup, including the microwave cavity in a superconducting magnet and the associated cryogenic and room-temperature receiver chain. }
    \label{fig:schematics}
\end{figure}


In the original search, near-quantum-limited sensitivity was achieved using multiple JPAs operated in a parallel configuration, in which a single pump tone simultaneously drove individual amplifiers with different frequency bandwidths.
The effective resonant frequency was selected by the pump frequency, enabling an extended search bandwidth while sharing a common circulator and receiver chain.
Although JPAs provided near-quantum-limited amplification over most of the frequency range, frequency intervals outside their bandwidth were necessarily scanned with HEMTs.
In one such dataset spanning 1.033--1.037\,GHz, the antenna coupling was not recorded, and the data were therefore excluded from the previous analysis.
We subsequently recovered the coupling by exploiting the smooth frequency dependence of the loaded quality factor $Q_L$, which was measured across the full range.
Motivated by this behavior, we assumed the antenna coupling to vary smoothly with frequency and inferred $\beta$ in the missing region by interpolating coupling values measured in adjacent frequency regions.
A quadratic interpolation was employed to reflect the overall frequency dependence.
The resulting reconstruction showed negligible boundary discontinuities, and the inferred $\beta$ was insensitive to the interpolation order, with a maximum deviation of $\sim$0.37\%.

During the analysis of the recovered dataset, an excess was observed at 1.036315\,GHz ($m_a=4.285857\,\mu$eV) in data acquired on January 23, 2023.
The excess exhibited a local statistical significance of $5.1\sigma$, which is reduced to a global significance of $3.5\sigma$ after accounting for the look-elsewhere effect.
The observed spectral profile was consistent with that expected from virialized axion dark matter~\cite{PhysRevD.42.3572}.
The corresponding power spectra are shown in Fig.~\ref{fig:signal_in_bins}.
The signal strength as a function of frequency tuning is consistent with the cavity resonance response, as shown in Fig.~\ref{fig:cavity_lineshape}.
The inferred signal strength corresponded to roughly $1.3$ times the KSVZ axion--photon coupling, assuming a local axion density of $0.45\,\mathrm{GeV/cm^{3}}$~\cite{10.1093/mnras/stae034, Lim_2025}.
Given its statistical significance and characteristic spectral profile, the excess was treated as an axion candidate and subjected to follow-up investigations.

\begin{figure}[ht]
    \centering
    \includegraphics[width=\linewidth]{figure/candidate_lineshape.pdf}
    \caption{Power spectra centered on the candidate frequency for different frequency-bin widths: (a) 50\,Hz and (b) 200\,Hz.
    The expected signal lineshape from the standard axion halo model is overlaid as red dashed lines.}
    \label{fig:signal_in_bins}
\end{figure}

\begin{figure}[ht]
    \centering
    \includegraphics[width=\linewidth]{figure/candidate_lorentzian.pdf}
    \caption{Signal strength as a function of frequency detuning $\nu-\nu_c$, normalized to the cavity bandwidth ($\Delta\nu_c=28.25$\,kHz).
    The error bars indicate statistical uncertainties.
    The data are fitted with a Lorentzian profile whose linewidth is consistent with the measured cavity resonance.}
    \label{fig:cavity_lineshape}
\end{figure}


A cross-check was performed using an independent haloscope system, previously utilized for an early axion search~\cite{CAPP-1}.
The apparatus consists of an 8-T superconducting magnet with a 165-mm bore and a cryogen-free dilution refrigerator operating below 40\,mK.
For this study, a dedicated cylindrical copper cavity with internal dimensions of $\O 134\,{\rm mm} \times 246\,{\rm mm}$ was fabricated.
Its natural TM$_{010}$ resonance at 1.713\,GHz was tuned down to the candidate frequency by inserting a $\O 30$-mm alumina rod.
This modification resulted in a reduced form factor of $C_{010} \approx 0.12$, as estimated from electromagnetic numerical eigenmode computations.
Quantum-limited amplification was achieved using a single flux-driven JPA optimized for the target frequency.
An initial dataset acquired in December 2023, taken with a cryogenic HEMT amplifier exhibiting elevated system noise ($T_{\rm sys}\sim500$\,mK), showed a transient excess at the candidate frequency with an observed SNR of $\sim$3.7, prompting further integration to assess its persistence; however, the excess did not reappear with increased exposure.
Data collected in a subsequent run in January 2024, after replacement of a cryogenic HEMT amplifier that reduced the system noise to $\sim$200\,mK, revealed no statistically significant excess, failing to confirm the candidate.

Following the independent cross-check, we re-examined the candidate using the original apparatus reconfigured with the same single JPA used in the cross-check experiment.
This configuration enabled a direct test of whether the previously observed excess originated from systematic effects associated with the amplification chain or other setup-specific artifacts.
A follow-up scan conducted between April and June 2024, comprising several segmented data-taking periods due to system instability that required repeated helium recondensing procedures, found no persistent signal with a spectral profile or temporal stability consistent with virialized axion dark matter.
We therefore find no evidence supporting a dark-matter interpretation of the candidate and conclude that the initially observed excess does not indicate galactic axion dark matter, although its precise origin remains unidentified.
We note that alternative scenarios involving clustered or transient axion dark matter have been proposed in the literature~\cite{Tinyakov_2016, PhysRevD.101.023008}; however, the present data are not sufficient to support such interpretations.

\begin{figure}[ht]
    \centering
    \includegraphics[width=0.95\linewidth]{figure/noise.png}
    \caption{Measured system noise temperature $T_{\rm sys}$ for this scan (blue), compared with the original scans employing JPA-based (green) and HEMT-based (yellow) receiver chains~\cite{CAPP-MAX}.
    The red dashed line indicates the standard quantum limit.}
    \label{fig:noise}
\end{figure}


To further strengthen this conclusion, we performed a targeted rescan over the $\sim$20-MHz band surrounding the candidate frequency.
Data was collected over a total of 12 days from June 20 to July 10, 2024, with an interruption due to facility maintenance.
The analysis procedure followed Ref.~\cite{CAPP-MAX}, consisting of baseline estimation, spectral preprocessing, and statistical combination.
Baselines extracted from cavity transmission and reflection spectra, JPA noise and gain spectra are used to determine the cavity parameters, expected signal power, and system noise temperature.
The performance of the baseline removal and normalization was verified by the resulting spectra following standard normal distributions.
Correlations between mirrored frequency bins induced by the JPA were accounted for by optimally combining correlated pairs using their expected SNRs and measured correlation coefficients, yielding an average SNR improvement of 4.1\%.
The analysis efficiency was quantified by injecting software-synthesized axion signals with a reference SNR of 5 directly into the obtained power spectra, and comparing the injected and recovered SNRs after the full analysis chain.
Repeating this procedure across the scanned frequency range, we found an average efficiency of $92.7\pm0.9\%$.

Spectra from different cavity tunings were combined using SNR weighting.
In the resulting spectrum, 146 clusters exceeding $5 \sigma$ were identified; all were found to be inconsistent with axion dark matter signals based on their narrow bandwidths, correlations with aerial antenna data, lack of the expected cavity response, or non-persistence upon rescanning, and were therefore removed together with neighboring bins.
Frequency bins were further combined within 5-kHz windows using the expected axion signal shape from the standard halo model.
A total of 16 clusters exceeding the rescan threshold of $3.718\sigma$ underwent follow-up measurements, yielding one persistent candidate, which was subsequently excluded as a galactic axion signal based on its persistence with the magnetic field turned off.
While incompatible with axion dark matter, such behavior is consistent with magnetic-field–independent scenarios, such as dark photon dark matter, and may warrant further investigation.

The dominant source of uncertainty in the analysis arises from the estimation of the system noise.
The noise temperature was measured at each frequency-tuning step, and its slow variation across the scanned frequency range was removed using a Savitzky–Golay filter~\cite{SGfilter}.
The residual fluctuations about the mean were taken as the associated uncertainty, yielding a fractional uncertainty of 6.4\%.
Figure~\ref{fig:noise} compares the system noise measured over the rescanned frequency range with that obtained in the original scan over the full frequency range.
Other analysis uncertainties, including those associated with the cavity parameters, contribute at the sub-1\% level.

With no signal consistent with axion dark matter observed, we set 90\% confidence-level upper limits on the axion--photon coupling over a frequency (mass) range 1.026--1.045\,GHz (4.244--4.322\,$\mu$eV), as shown in Fig.~\ref{fig:exclusion}.
Due to the performance of the JPA, the upper portion of this range achieved sensitivity approaching the DFSZ benchmark, while the lower portion reached sensitivity at the KSVZ level.
These results supersede previous constraints and represent among the most stringent haloscope limits reported to date in this range.

In summary, we completed an axion haloscope search by recovering the previously unrecorded frequency region near 1.036\,GHz from Ref.~\cite{CAPP-MAX}.
A candidate excess was identified and subjected to a rigorous program of validation, ultimately failing to confirm as an axion dark matter signal.
Independent measurements by the ADMX experiment also probed this frequency region with adequate sensitivity and found no indication of an excess signal (G. Rybka, private communication).
We subsequently extended search around the candidate frequency using a quantum-noise-limited JPA receiver, enabling improved 90\% confidence-level upper limits on the axion–photon coupling.
This work demonstrates the robustness and importance of systematic candidate-validation strategies as axion haloscope searches approach discovery-level sensitivity.

\begin{comment}
\begin{table}
    \centering
    \caption{Uncertainties quoted in the analysis, expressed as percentages relative to the central values of the corresponding parameters.}
    \begin{tabular*}{\columnwidth}{@{\extracolsep{\fill}}ccc}
        \toprule
        Parameter        & Uncertainty [\%] &  Source \\
        \midrule
        $\epsilon_{\rm SNR}$ & 1.0 & Frequency-dependent variation \\
        $\beta$ & 3.5 & Fitting uncertainty\\
        $T_{\rm sys}$ & 6.2 & Statistical error in $G^{\rm JPA}$ and $T_{\rm off}$\\
        $Q_L$ & $<$ 1.0 & Statistical fluctuation \\
        $C$ & $<$ 1.0 & Thermal contraction effects \\
        \bottomrule
    \end{tabular*}
    \label{tab:uncertainty}
\end{table}
\end{comment}


\begin{figure*}[ht]
    \centering
    \includegraphics[width=\linewidth]{figure/exclusion.png}
    \caption{Experimental limits on the axion–photon coupling $g_{a\gamma\gamma}$ as a function of axion mass $m_a$ (frequency $\nu_{a}$).
    The upper panel shows the 90\% confidence-level limits set by this work, assuming a local axion density of $\rho_a=0.45\,{\rm GeV/cm^3}$; the black dashed lines indicate the KSVZ and DFSZ benchmark models.
    The vertical red dashed line denotes the candidate frequency.
    The lower panel compares the results of this work with previously reported haloscope limits~\cite{ADMX, ADMX2018, ADMX2019_1_2, ADMX2021, ADMX2024, ADMX2025, ADMX_Sidecar, ADMX_Sidecar_JTWPA, CAPP-1, CAPP-2, CAPP-3, CAPP-4, CAPP-5, CAPP-6, CAPP-7, CAPP-8, CAPP-9, CAPP-MAX, CAPP-8TB-PR2, DMAG-12T, CASTCAPP, GrAHal, HAYSTAC2, HAYSTAC3, HAYSTAC4, ORGAN, QUAX, QUAX2, QUAX3, QUAX4, QUAX5, RADES, RADES2, RBF, UF, TASEH}, with the exclusion data taken from the AxionLimits GitHub repository~\cite{AxionLimits}.}
    \label{fig:exclusion}
\end{figure*}


\begin{acknowledgments}
This work was supported by the Institute for Basic Science (IBS-R017-D1-2024-a00 and IBS-R040-C1-2025-a00) and by JSPS KAKENHI (Grant No.~JP22H04937).
A. F. van Loo was supported by a JSPS Postdoctoral Fellowship.
J. E. Kim was partially supported by the Korea Academy of Science.
\end{acknowledgments}



\bibliography{main}

\end{document}
===== END FILE: main.tex =====
